% !TEX TS-program = pdflatex
% Overleaf: blank project is fine. Menu -> Compiler = pdfLaTeX.
% No extra files. Theme Madrid is in TeX Live.
\documentclass[aspectratio=169]{beamer}
\usetheme{Madrid}
\usecolortheme{default}
\usepackage{booktabs,array}
\setbeamertemplate{navigation symbols}{}
\setbeamertemplate{footline}[frame number]

\title{The Judger already writes as fast as this GPU allows}
\subtitle{Hybrid LatentMAS serving $\cdot$ Qwen3-14B $\cdot$ one H200 $\cdot$ greedy}
\author{Sagnik Biswas}
\date{13 October 2026}

\begin{document}

\begin{frame}
\titlepage
\end{frame}

\begin{frame}{What this round was for}

The earlier result still stands: on a long answer, the three silent agents are
about a second, and the Judger's written answer is the rest of the clock.

\vspace{0.6em}
This round asked one question. \textbf{Can the serving stack make that clock
shorter}, with the model, the prompts, the latent-step count, and the decoding
policy held fixed?

\vspace{0.6em}
\begin{itemize}
  \item Starting point: the authors' optional fast path. Planner, Critic, and
        Refiner stay in Hugging Face. The Judger runs in vLLM.
  \item Hardware: one H200. Stack: PyTorch 2.8, Transformers 4.57.1, vLLM 0.11.0.
  \item The authors' published numbers are the Hugging Face-only path. This
        hybrid path is a separate implementation. It is not claimed to match
        those numbers bit for bit.
\end{itemize}

\end{frame}

\begin{frame}{Where one request spends its time}

Repaired hybrid path, one question, 1024 forced tokens, greedy.

\begin{center}
\begin{tabular}{lcc}
\toprule
piece & seconds & share \\
\midrule
three latent agents, 10 steps each & 1.07 & 10\% \\
hand the vectors to vLLM & $\sim$0 & --- \\
Judger writes 1024 tokens & 9.37 & 89\% \\
\midrule
end to end & 10.49 & 100\% \\
\bottomrule
\end{tabular}
\end{center}

\vspace{0.5em}
\begin{itemize}
  \item The latent shortcut is doing its job. Those agents never sample a word.
        They feed the hidden vector back in as the next input.
  \item The Judger cannot skip that. A person needs the written answer.
  \item Moving the embedding tensor onto the GPU is not the cost.
\end{itemize}

\end{frame}

\begin{frame}{The authors' hybrid code has two real bugs}

These are correctness bugs. Fixing them does not make a request faster.
The repaired path is what every timing below uses. The original path is
still the default, so the crash stays reproducible.

\vspace{0.4em}
\begin{itemize}
  \item \textbf{Batch size 1 crashes.} The code squeezes a 3D embedding down
        to 2D, then tries to unpack three dimensions.
        \texttt{ValueError: not enough values to unpack (expected 3, got 2)}.
        Recorded after a full model load.
  \item \textbf{A batched call is not the same as separate calls.} Padding
        positions are kept, and the pad value is float32, so a bf16 embedding
        is promoted. The repaired path drops masked positions and keeps bf16.
\end{itemize}

\end{frame}

\begin{frame}{Giving vLLM vectors does not slow the writing}

vLLM normally wants token ids. The hybrid Judger has to hand it embedding
rows instead. vLLM's own comments say that mode leaves the embedding lookup
outside the compiled graph, which sounded like a tax on every generated token.

\vspace{0.4em}
Same engine, same GPU, 1024 greedy tokens, no latent agents in either run.

\begin{center}
\begin{tabular}{lcc}
\toprule
how the engine was built & seconds & tokens/s \\
\midrule
ordinary text prompt & 9.37 & 109.3 \\
prompt embeddings enabled & 9.51 & 107.7 \\
\bottomrule
\end{tabular}
\end{center}

\vspace{0.4em}
A raw memory copy on this H200 hits 4.26 TB/s. A 14B bf16 weight read then
caps batch-1 decode near 152 tokens/s. \textbf{109 is ordinary batch-1
decode}, about 72\% of that ceiling.

\end{frame}

\begin{frame}{Hiding the next latent rollout behind the Judger does not pay}

On one GPU the two pieces contend. Four requests, 1024 forced tokens.

\begin{center}
\begin{tabular}{lc}
\toprule
arm & wall clock \\
\midrule
one request fully finishes, then the next starts & 41.52 s \\
build request $i{+}1$ while vLLM writes request $i$ & 40.75 s \\
\bottomrule
\end{tabular}
\end{center}

\vspace{0.4em}
\begin{itemize}
  \item Speedup \textbf{1.019$\times$}.
  \item Decode slowed from about 9.6 s to about 10.1 s whenever the next
        latent rollout ran on the same card. That cancels the $\sim$0.8 s
        of prep.
  \item The greedy text matched across arms, so the overlap did not corrupt
        the answer. This pod had one GPU. The authors' two-GPU split was
        not measured.
\end{itemize}

\end{frame}

\begin{frame}{Prefix caching does not engage}

vLLM 0.11 turns prefix caching off when prompt embeddings are enabled.
We checked the live engine, not a timing guess.

\begin{itemize}
  \item \texttt{enable\_prefix\_caching} on the running engine: \textbf{false}.
  \item Replaying the identical embedding, and a copy with the middle row
        shifted by 1: \texttt{num\_cached\_tokens} $= 0$ both times
        (0.097 s vs 0.095 s).
  \item The first eight tokens matched because both answers open with
        \texttt{<think>}. That is the model, not a cache hit.
\end{itemize}

\end{frame}

\begin{frame}{On real questions the number is right, and the essay is the cost}

Three saved questions. Greedy, stop on EOS, cap 512. None hit the cap.

\begin{center}
\begin{tabular}{lrrrr}
\toprule
question & tokens & end to end & of which vLLM & answer \\
\midrule
add    & 257 & 3.48 s & 2.40 s & 42 \\
packs  & 293 & 3.56 s & 2.73 s & 30 \\
long   & 499 & 5.82 s & 4.86 s & 37 \\
\bottomrule
\end{tabular}
\end{center}

\vspace{0.4em}
\begin{itemize}
  \item All three matched the gold integer.
  \item Each answer is a long \texttt{<think>} trace that ends in
        \texttt{\textbackslash boxed\{42\}} (or 30, or 37).
  \item Latent rollout stays about 0.8--1.0 s. Once the answer is a few
        hundred tokens, vLLM is most of the request.
\end{itemize}

\end{frame}

\begin{frame}{Batching is the throughput lever that already works}

Fixed 1024 tokens, so a shorter answer cannot fake a speedup.
One wall-clock number per batch.

\begin{center}
\begin{tabular}{lccc}
\toprule
batch & wall & latent & vLLM \\
\midrule
1 & 10.49 s & 1.07 s & 9.37 s \\
2 & 10.36 s & 0.85 s & 9.50 s \\
4 & 10.54 s & 0.89 s & 9.63 s \\
8 & 13.00 s & 2.25 s & 10.68 s \\
\bottomrule
\end{tabular}
\end{center}

\vspace{0.4em}
Eight requests finish in 13.0 s. Eight serial requests would be about 84 s.
That is about \textbf{6.5$\times$} the single-request rate, from the offline
batch the code already had. Zero errors across this sweep.

\end{frame}

\begin{frame}{What to take from this}

\textbf{We did not make a single request faster than the hybrid path we
started from.}

\vspace{0.5em}
\begin{itemize}
  \item The Judger writes at $\sim$109 tokens/s. A plain text prompt on the
        same vLLM build writes at the same speed.
  \item Previously, the Hugging Face Judger on this model measured
        $\sim$42.5 tokens/s. vLLM was already that gain. It is a different
        workload, so it is context, not a paired A/B from this week.
  \item The repair is worth keeping: batch size 1 no longer crashes, and a
        batch matches the requests that make it up.
  \item The open lever is how many words the Judger writes. Same gold
        number, fewer tokens, and the wall clock falls with them. A shorter
        wrong answer does not count.
\end{itemize}

\end{frame}

\end{document}
